\section{Testing at a prescribed Hamming distance}\label{sec:test}
A tester must accept $x\in A_s$ and reject tables with $d(x,A_s)\ge\xi N$, where $d$ counts differing entries. Its output elsewhere is unrestricted. Let $R_\xi,Q_\xi$ be the classical randomized and quantum query costs with pointwise error at most $1/3$ on these two cases. Computation between queries is free.
\begin{theorem}\label{th:testing}
Fix $\eta,\epsilon>0$. For some $\xi_0>0$ and every $1/N\le\xi\le\xi_0$, in the gate range of Theorem~\ref{th:main}, write $L_\xi=1+\log(1/\xi)$ and $K_N=T/(1+\log(N/T))$. Then
\begin{align*}
 R_\xi &=\Theta_{\eta,\epsilon}\!\left(\min\{N,\, T/(\xi L_\xi)\}\right),\\
 Q_\xi &=\Theta_{\eta,\epsilon}\!\left(\min\{\sqrt{NK_N},\,T/(\sqrt\xi L_\xi)\}\right).
\end{align*}
\end{theorem}
Thus at constant proximity in this range both models use order $s\log s$ queries. As the allowed distance decreases, knowing the input is far on a no-instance still helps until $\xi$ reaches order $K_N/N$.

\paragraph{An upper bound for any finite family.}
Let $A$ be a known class of at least two binary strings of length $N$ and let $B=\log |A|$. Choose $a=\xi$ classically and $a=\sqrt\xi$ quantumly. Searching for a coordinate at which two tables differ can be done with a constant probability of finding a mismatch using $O(1/a)$ queries whenever they disagree on at least $\xi N/3$ coordinates. Use random sampling or quantum search with the given lower bound on the marked fraction, always verifying any proposed coordinate. Set $\gamma=c a\log(2/a)$ with a sufficiently small absolute $c$ so $\gamma<1/2$.

Keep only candidates consistent with actual queried values. If a coordinate splits the survivors with minority proportion at least $\gamma$, query it. At most $O(B/\gamma)$ such queries occur since each retains at most a $(1-\gamma)$ fraction. Otherwise let $m$ be the coordinatewise majority of the survivors. Cast one vote for acceptance if some survivor lies within $\xi N/3$ of $m$ and for rejection otherwise. Then make one search attempt for a disagreement of $x$ with $m$. A found mismatch retains at most the $\gamma$ fraction. If an attempt finds nothing, the survivors stay in place. When they become empty reject; for a singleton verify its member by a single amplified promise search, distinguishing equality from distance $\xi N$.

Here is a uniform bound on the accumulated error. On a promised input, whenever the vote is wrong, $d(x,m)\ge\xi N/3$. Indeed if $x\in A$ it is still a survivor; if $x$ is far then the triangle inequality applies to the proposed nearby survivor. A fixed survivor set with a wrong vote therefore contributes at most a constant expected number of wrong votes before a mismatch is found. There are at most $1+B/\log(1/\gamma)$ distinct sets at which votes are cast: every successful minority search reduces the survivor count by a factor $\gamma$. Stop after a large constant times $1+B/\log(1/\gamma)$ votes or earlier termination. If the vote cap is reached, output the majority. Markov's inequality bounds the probability of a wrong majority, with any desired fixed error margin by choice of cap. Balanced queries and witnessed mismatches never remove the target on a yes-input. This proves the query bound
\[
 O\!\left(\frac1a+\frac{B}{a\log(2/a)}\right).
\]
Use (1), the exact upper bounds, and $T\gg\log N$ to obtain both claimed upper bounds.

\paragraph{Separation from a Hamming neighborhood.}
For an ambient length $M=2^m$ under Theorem~\ref{th:main}'s gate hypotheses, set $K_M=T/(1+\log(M/T))$. There is an $\alpha>0$, depending only on the fixed endpoint margins, such that the two experiments in our proof separate $A_s$ from distance at least $\alpha K_M$ using lower bounds $\Omega(M)$ classical queries and $\Omega(\sqrt{M K_M})$ quantum queries. For the distance assertion, enlarge the easy family in the hard-law union bound to its Hamming neighborhood. In the superlinear sparse range a radius $\alpha K_M$ adds at most $(O(\alpha)+o(1))T$ bits to the log count. This follows from $\sum_{j\le r}\binom Mj\le (eM/r)^r$ for $1\le r\le M/2$, and $\log(M/K_M)\sim1+\log(M/T)$ when $T=o(M)$. The hard point-probability exponent has a fixed fractional margin over $T$. The same reasoning applies at bounded linear gate ratios, taking the number of copies large enough that the exponent in Section~4 exceeds $2T$. At a positive Shannon fraction the uniform hard law gives the neighborhood probability bound $|A_s|2^{-M}\sum_{j\le\alpha M}\binom Mj=o(1)$ if binary entropy $H_2(\alpha)<\epsilon/2$. The constants can be taken uniformly by the subsequence argument of Section~4.

For the classical query lower bound, fix a deterministic adaptive $u$-query algorithm and fix the shifts. Each query to (2) either reads a known mask bit or reveals a fresh independent Bernoulli bit. Its expected output at complex parameter $z$ is a polynomial of degree at most $u$ and modulus bounded by $(|z|+|1-z|)^u$. To see this, write it as a sum of transcript weights, retaining the value of a bit after its first read, and sum their absolute values down the binary decision tree. Average over algorithm coins and over the two label laws. Their expectation difference $R$ has $|R(z)|\le2\exp(2u r)$ on $|z|\le r$. The decision polynomial is multilinear of degree at most $u$, so the same touched-row argument as in Lemma~\ref{lem:analytic} gives a zero of order $D$ at zero. At $r=4t_0$, $|R(t_0)|\le 2\cdot4^{-D}\exp(8u t_0)$, while the promise forces a constant gap in the real acceptance probabilities. Thus $u=\Omega(DH)=\Omega(M)$ in the mask experiments. In the dense code experiment, bounded independence already fools every deterministic decision tree of depth at most $\delta M$, and hence their randomized mixtures. For the quantum lower bound in all cases, the decision algorithm gives an acceptance polynomial bounded throughout the Boolean cube, even on inputs outside its promise.

\paragraph{Adapting the domain length.}
Fix $c_0>0$ less than the neighborhood constant $\alpha$, using a Shannon margin of $\min(\epsilon,1/2)$ in that constant. Choose $\xi_0$ sufficiently small in terms of $c_0$. Let $M$ be the greatest power of two at most
\[
                 \min\{N,\ c_0 T/(\xi L_\xi)\}.
\]
If $M<N$, we can ensure $M\ge4T$ and $1+\log(M/T)=\Theta(L_\xi)$ by shrinking $\xi_0$. Since $m/2^m$ is decreasing for $m\ge2$, $s m/M\le1/2$ for large $s$, by checking near $M=4T$. If $M=N$, the original gate margins apply. In both cases $1+\log(M/T)\le L_\xi$, whence $\xi M\le\alpha K_M$. Duplicate a table on $m$ address bits across all assignments to the remaining $n-m$ bits. Its minimum circuit size is unchanged by duplication since any ambient circuit can be restricted. Each such restriction of a size-$s$ circuit must disagree with a hard $m$-variable table on at least $\alpha K_M$ positions. Summing these discrepancies over the slices shows the padded table is $\xi N$-far. One query to it costs one query to the original table, coherently as well as classically. For $M=N$ this recovers the exact lower scale. For $M<N$, substitute $M=\Theta(T/(\xi L_\xi))$ and $K_M=\Theta(T/L_\xi)$ into the preceding separation, completing the proof.

The testing-by-implicit-learning literature gives classical constant-distance upper bounds via finite families on relevant inputs. Alon, Blais, Chakraborty, Garc\'ia-Soriano and Matsliah prove a polynomial-in-$s$ query lower bound for testing the entire class of size-$s$ circuits at constant distance (their Section~7.4, Corollary~2.2). Their work also studies function isomorphism \cite{iso,extractors}. Our calculation gives both query models at varying proximity for fixed binary fan-in and the gate range in Theorem~\ref{th:testing}. It makes no efficiency assumption on enumerating candidate circuits.
